% ============================================================================
%  SECCIÓN 3.24: COMPARACIÓN WIREFRAME VS IMPLEMENTACIÓN
% ============================================================================

\section{Comparaci\'on wireframe vs implementaci\'on}

Siguiendo la gu\'ia del avance, se realiza una comparaci\'on entre los
wireframes dise\~nados en el primer avance (secci\'on de wireframes iniciales) y
las pantallas implementadas en esta etapa. La tabla resume los elementos
conservados, modificados, agregados y pendientes en cada pantalla.

\begin{table}[H]
  \centering
  \caption{Comparativa entre wireframes e implementaci\'on}
  \label{tab:comparacion-wireframe}
  \contenidoConFuente{%
    \begin{tablaAPA}
      \begin{tabular}{@{}p{2.6cm}p{3.0cm}p{3.2cm}p{3.0cm}p{3.2cm}@{}}
        \toprule
        \textbf{Pantalla} & \textbf{Conservados} & \textbf{Modificados} &
        \textbf{Agregados} & \textbf{Pendientes} \\\
        \midrule
        Splash & Logo, nombre de la app, bot\'on de acci\'on &
        Fondo con forma decorativa curva en lugar de fondo plano &
        Mensaje tagline \textit{``Descubre, explora y conecta''} &
        Animaci\'on de transici\'on \\\
        \addlinespace
        Login & Campos de email y contrase\~na, bot\'on de acci\'on, enlace a
        registro & Layout centrado vertical con espaciado amplio &
        Indicador de carga, mensajes de error de validaci\'on &
        Iconos en los campos de texto \\\
        \addlinespace
        Registro & Campos de nombre, email, contrase\~na, bot\'on de registro &
        Agregado campo de confirmaci\'on de contrase\~na &
        Validaci\'on en tiempo real con Zod, mensaje de \'exito,
        redireccionamiento autom\'atico &
        Selector de rol de usuario \\\
        \addlinespace
        Principal & Barra de navegaci\'on inferior con 4 secciones &
        Las 4 pesta\~nas reemplazan el dise\~no de mapa propuesto &
        \'Iconos de Ionicons en cada pesta\~na, colores del tema &
        Mapa interactivo, barra de b\'usqueda \\\
        \addlinespace
        Detalle & Nombre del escenario, imagen, informaci\'on b\'asica &
        Recibe ID por ruta din\'amica &
        Placeholder funcional que valida la navegaci\'on &
        Imagen del escenario, disciplinas, horarios, bot\'on de favorito \\\
        \bottomrule
      \end{tabular}
    \end{tablaAPA}%
  }{Elaboraci\'on propia en base a la comparaci\'on entre los wireframes del
  primer avance y la implementaci\'on actual.}
\end{table}

\subsection{An\'alisis por pantalla}

\subsubsection{Pantalla Splash}

El wireframe original mostraba un dise\~no simple con logo, nombre y bot\'on.
La implementaci\'on conserva estos tres elementos y agrega un fondo decorativo
con forma curva que genera una identidad visual m\'as fuerte. El mensaje tagline
se agrad\'o como elemento adicional para comunicar la propuesta de valor de la
aplicaci\'on.

La implementaci\'on actual de la pantalla Splash se presenta en la
figura~\ref{fig:pantalla-splash} de la secci\'on de construcci\'on de
pantallas. El wireframe original del primer avance se presenta en la
figura~\ref{fig:wireframes-acceso} de la secci\'on de wireframes iniciales.
La comparaci\'on muestra que se conservaron los elementos fundamentales
(logo, nombre, bot\'on de acci\'on) y se a\~nadi\'o un fondo decorativo
y un mensaje tagline para reforzar la identidad visual.

\subsubsection{Pantallas de login y registro}

Los wireframes originales propon\'ian campos de texto simples con un bot\'on de
acci\'on. La implementaci\'on conserva esta estructura y agrega
validaci\'on en tiempo real, mensajes de error, indicadores de carga y
enlaces de navegaci\'on entre pantallas. El formulario de registro incorpora
un campo adicional de confirmaci\'on de contrase\~na que no exist\'ia en el
wireframe original, mejorando la seguridad de la experiencia de registro.

\subsubsection{Pantalla principal}

El wireframe original propon\'ia una pantalla con barra de b\'usqueda y mapa
interactivo. La implementaci\'on actual presenta una barra de navegaci\'on
inferior con cuatro pesta\~nas, lo que proporciona una estructura de
navegaci\'on m\'as clara. El mapa interactivo y la barra de b\'usqueda est\'an
pendientes de implementaci\'on en avances posteriores.

\subsubsection{Pantalla de detalle}

El wireframe propon\'ia una ficha completa con imagen, disciplinas, capacidad,
horarios y bot\'on de favorito. La implementaci\'on actual es un placeholder
funcional que valida la navegaci\'on y la recepci\'on de par\'ametros por ruta
din\'amica, dejando el dise\~no visual completo para un avance futuro.
